\documentclass[../../../include/open-logic-section]{subfiles}

\begin{document}

\olfileid{sth}{story}{blv}

\olsection{پیوست: قانون اساسی پنجم فرگه}

در~\olref[rus]{sec} توضیح دادیم که راسل پارادوکس خود را به‌صورت مشکلی
برای دستگاهی صورت‌بندی کرد که فرگه در~\emph{Grundgesetze} ترسیم کرده بود.
دستگاه فرگه صورت‌بندی مستقیمی از تصریح ساده‌انگارانه نداشت. پس در این
پیوست بسیار کوتاه توضیح خواهیم داد که دستگاه فرگه چه چیزی را
\emph{دربر می‌گرفت}، و چه نسبتی با تصریح ساده‌انگارانه و پارادوکس راسل
داشت.

دستگاه فرگه \emph{مرتبهٔ دوم} بود و برای صورت‌بندی مفهوم
\emph{مصداق یک مفهوم} طراحی شده بود.\footnote{به‌بیان دقیق، فرگه می‌کوشد
مفهومی عام‌تر را صوری کند: «دامنهٔ مقادیر» یک تابع. مصداق‌های مفاهیم حالت
خاصی از این مفهوم عام‌ترند. برای جزئیات بنگرید به
\citet[pp.\ 8--9]{Heck2012}.} با نمادگذاری‌ای الهام‌گرفته از فرگه،
$\fregeext{x}{F(x)}$ را برای \emph{مصداق مفهوم~$F$} می‌نویسیم. این
سازوکاری است که یک \emph{محمول}، یعنی «$F$»، را می‌گیرد و آن را به یک
\emph{ترم} (مرتبهٔ اول)، یعنی «$\fregeext{x}{F(x)}$»، تبدیل می‌کند. فرگه
با استفاده از این سازوکار، تعریف زیر را برای عضویت عرضه کرد:
\[
	a \in b \liff_{\text{df}} \exists G(b = \fregeext{x}{G(x)} \land Ga)
\]
یعنی به‌تقریب، $a \in b$ اگر و تنها اگر~$a$ ذیل مفهومی قرار گیرد که
مصداق آن~$b$ است. (توجه کنید سور «$\exists G$» مرتبهٔ دوم است.) فرگه اصل
زیر را نیز می‌پذیرفت که به \emph{قانون اساسی پنجم} مشهور است:
$$\fregeext{x}{F(x)} = \fregeext{x}{G(x)} \liff \forall x (Fx \liff Gx)$$
یعنی به‌تقریب، مفاهیم مصداق‌های یکسان دارند اگر و تنها اگر هم‌مصداق باشند.
(باز هم، هر دو «$F$» و «$G$» در جایگاه محمول قرار دارند.) اکنون اصلی ساده
عضویت را به صدق یک ویژگی دربارهٔ شیء پیوند می‌دهد:

\begin{lem}[در \emph{Grundgesetze}]\ollabel{lem:Fregeextensions}
$\forall F \forall a(a \in \fregeext{x}{F(x)} \liff Fa)$
\end{lem}

\begin{proof}
$F$ و~$a$ را ثابت بگیرید. آنگاه
$a \in \fregeext{x}{F(x)}$ اگر و تنها اگر
$\exists G(\fregeext{x}{F(x)} = \fregeext{x}{G(x)} \land Ga)$
(بنا بر تعریف عضویت)، اگر و تنها اگر
$\exists G(\forall x(Fx \liff Gx) \land Ga)$
(بنا بر قانون اساسی پنجم)، اگر و تنها اگر~$Fa$ (بنا بر منطق مقدماتی
مرتبهٔ دوم).
\end{proof}

و از این نتیجه تقریباً بی‌درنگ تصریح ساده‌انگارانه به دست می‌آید:

\begin{lem}[در \emph{Grundgesetze}.]
$\forall F \exists s \forall a (a \in s \liff Fa)$
\end{lem}

\begin{proof}
$F$ را ثابت بگیرید؛ اکنون~\olref{lem:Fregeextensions} نتیجه می‌دهد
$\forall a (a \in \fregeext{x}{F(x)} \liff Fa)$؛ پس با تعمیم وجودی
$\exists s\forall a(a \in s \liff Fa)$. چون~$F$ دلخواه بود، نتیجه به
دست می‌آید.
\end{proof}

پارادوکس راسل با گرفتن~$F$ به‌گونه‌ای که
$\forall x(Fx \liff x \notin x)$ باشد نتیجه می‌شود.

\end{document}
